\section{Conclusiones}
\label{sec:conclusiones}

La segmentación cumplió su objetivo, con un Dice medio de 0,916 en prueba y fallas en sólo 4 de las 200 imágenes. La compresión quedó caracterizada, aunque no se llegó al objetivo de máxima de repetir el algoritmo sobre imágenes comprimidas.

La hipótesis 1 no se confirmó. Quitar los pelos con un cierre no mejoró la segmentación y, en entrenamiento, resultó peor que el filtro de mediana. Una posible razón, que no comprobamos, es que el suavizado gaussiano previo a la segmentación ya atenúa buena parte de los pelos.

El criterio D estimado con un solo atributo de textura no reproduce al dermatólogo y, en prueba, tampoco separa los grupos con seguridad. El AUC cayó de 0,73 en entrenamiento a 0,60 en prueba, y eso muestra por qué convenía reservar ese conjunto. Mirar la lesión por zonas fue más prometedor, porque la saturación de periferia menos centro se mantuvo en un AUC de 0,71, pero con 60 imágenes el intervalo es ancho y no alcanza para distinguirla del brillo de la lesión ni para afirmar que mejore a D.

A, B y C se pudieron medir, aunque sólo en parte como el dermatólogo. A coincide con él de forma aceptable ($\kappa=0{,}38$) y falla cuando la asimetría está en el color y no en la forma, C coincide poco ($\kappa=0{,}24$) y B separa poco porque en PH2 los nevos comunes tienen el borde más nítido. La circularidad, que no es parte de la regla, separó casi igual que A, y es una buena candidata a término extra del TDS.

La hipótesis 2 se confirmó sólo en parte. El TDS automático separa las lesiones sospechosas de los nevos comunes (AUC 0,78), pero agregarle las zonas no lo mejora. Con el umbral de Youden su sensibilidad es 0,78 y su especificidad 0,58, lejos del 91\,\% y el 96\,\% de \textcite{kasmi2016}, aunque ese trabajo separa melanomas de lesiones benignas, y acá los nevos atípicos cuentan como sospechosos y su TDS se parece más al de los nevos comunes. Con el umbral clínico, el programa reconoce 8 de los 12 melanomas de prueba.

Las limitaciones principales son la base, chica y de un solo equipo y un solo hospital, y el conjunto de prueba, que tiene sólo 60 imágenes y 9 con pelos. Además, D es una adaptación que cuenta estructuras con un número global, la referencia del rojo sale de una sola zona y la segmentación sobre el canal azul pierde justamente las zonas azul grisáceas que se asocian al melanoma, aunque con la máscara manual los criterios no mejoraron.

Como mejoras, proponemos segmentar con información de color sin sumar fallas en el resto de las imágenes, medir cada estructura de D con su propio detector en lugar de un único número, repetir el algoritmo sobre imágenes comprimidas y validarlo con otra base de imágenes.
